\section{Conclusion}

\label{sec:conclusion}

We have identified temporal dominance as a structural failure mode in the deployed hybrid memory selector, derived the exact algebraic condition under which it occurs, and proposed CTLS as a principled correction. The deployed hybrid is outperformed by BM25 by -36.1111 percentage points, -29.3651 percentage points, and -33.7302 percentage points at the three evaluated token budgets in a controlled evaluation where only the selector changes. The temporal term causes most of this deficit: removing it (hybrid\_no\_time) recovers the majority of the gap. CTLS eliminates temporal dominance by construction through tier partition, preserving temporal decay as a within-tier tiebreaker.

The Non-Dominance Property proved in the method section guarantees that no zero-overlap item can outrank a positive-overlap item under CTLS, regardless of freshness gap, frequency, or importance values. This is a structural guarantee, not a hyperparameter constraint: it holds for any positive value of the temporal weight within the relevant-tier scoring function. CTLS requires no model calls, no embedding model, no learned parameters, and no write-path changes. It is a zero-cost selector substitution that can replace the deployed hybrid in any context that uses the same frozen-bank, whole-item-packing retrieval architecture.

The practical contribution of this paper is the combination of a formal diagnosis and a constructive fix. The diagnosis identifies the exact condition under which a common retrieval heuristic fails, shows why the failure is invisible when evaluated only against a weak baseline, and shows that a large margin against a strong lexical baseline plus a clean ablation together pinpoint the cause. The fix is provably correct, immediately actionable, and requires no new infrastructure. Future work should directly evaluate CTLS on the existing frozen bank under the same controlled conditions, complete the blinded judge adjudication, and extend the mechanism analysis to multi-session and temporal-reasoning question types.
